% Nouveau EGA V (cote 25): archive cote 25, site page 12 (PDF page 13).
% typescript. Draft compared: none.
% AI-transcribed by Claude Opus 5.5 (claude-opus-5-5, adaptive thinking, effort high), Message Batches API, batch ega25b (msgbatch_01A7n5YuXubUPNK1zsepi8RG), 2026-11-13; not yet checked by eye.
\scanpage{25/12}
\mspage{\hand{9}}

\uline{Corollaire} 2.11. Soit $y\in Y_{\hand{y}}$ , pour que $Y_{\hand{y}}$ \inserted{soit régulier (resp.} satisfasse la propriété
$R_k$ (réf) en y\hand{,} \deleted{(}resp. soit normal en y) il f et s que X satisfasse
la même propriété en x. \xout{\unsure{Remarque}} Soit Z l'ensemble des points de X en
lesquels X n'est pas \deleted{\xout{\illegible{1}} (}\inserted{régulier (resp. \deleted{$R_k$} $R_k$}resp. \deleted{n'est pas} normal); pour que $Y_{\hand{y}}$ soit
\inserted{régulier (resp.} $(R_k)$\hand{,} \deleted{(}resp. normal) il f et s que f(Z) soit fini i.e. dim $\overline{f(Z)}=0$.

\margin{|} Même démonstration que 2.8. et 2.9. Il faut donner les diverses
références assurant que Z est fermé (car il faut savoir qu'il est
constructible pour appliquer 2.3.).

\xout{\unsure{Remarque 2.12}} On notera que dans 2.10. et 2.11. nous ne disons rien
sur les propriétés ``géométriques'' correspondantes, les résultats
envisagés étant de nature \hand{``}absolue\hand{''}. Nous Examinons maintenant les
propriétés de nature géométrique. (On pourrait si on le désire en pro-%
\margin{||}fiter pour changer de n$^{\text{o}}$).

\inserted[typed]{Supposons que $f:X\to P$ soit une \deleted{\unsure{immersion}}\inserted{\unsure{non ramifié|non ramifiée}}.}
\uline{Théorème} 2.12. Soient $y\in Y_{\hand{y}}$ , x son image dans X. Pour que X soit
lisse sur k en x, il f et s que $Y_{\hand{y}}$ soit lisse sur $k(\hand{y})$ en y.

On peut supposer k algébriquement clos. Si $Y_{\hand{y}}$ est lisse sur $k(\hand{y})$
en y, il est régulier, donc comme $Y_{\hand{y}}$ estb plat sur X, X est régulier
en x (réf), donc il est lisse sur k en x puisque k est alg clos donc
parfait (réf). Pour la réciproque, on peut (quitte à remplacer X par
un voisinage ouvert de x) supposer X lisse, et (grâce au critère jaco-%
bien de lissité) défini dans un ouvert U de P par p équations
comme $X=V(f_1,..,f_p)$ , où les différentielles $df_1,..,df_p$ sont
partout linéairement indépendantes. Introduisant d es coordonnées affines
$S_1,..,S_n$ dans $\overset{\hand{\vee}}{P}$ et des coordonnées affines $T_1,..,T_n$ dans un voisinage
de x (par choix d'un hyperplan $H_\infty$ ne contenant pas x), $Y\hookrightarrow U_{k(\hand{y})}$ est
alors égal à $V(f_1,..,f_p,\hand{\sum_i}S_iT_i-1\ )$ , et il suffit de vérifier que
les différentielles (relativement à $k(\hand{y})$ ) de $f_1,..,f_p$ ,$\hand{(\sum}S_iT_i\hand{)}-1$

\mspage{\hand{12}}
